% --- Chapter: Variables and Types ---
\section{Variables and Types}

\subsection{Value Types}

CRISP supports dynamic typing with the following value types:

\begin{table}[H]
\centering
\rowcolors{2}{softviolet!30}{white}
\caption{\textcolor{darkpurple}{CRISP Value Types}}
\vspace{0.3em}
\begin{tabular}{@{}lllp{5.5cm}@{}}
\toprule
\rowcolor{softvioletdark}
\textbf{\textcolor{white}{Type}} & \textbf{\textcolor{white}{Rust Type}} & \textbf{\textcolor{white}{Example}} & \textbf{\textcolor{white}{Description}} \\
\midrule
Null & \texttt{Value::Null} & \texttt{null} & Absence of value \\
Bool & \texttt{Value::Bool} & \texttt{true}, \texttt{false} & Boolean values \\
Int & \texttt{Value::Int(i64)} & \texttt{42}, \texttt{-7} & 64-bit signed integer \\
Float & \texttt{Value::Float(f64)} & \texttt{3.14}, \texttt{-0.5} & 64-bit floating point \\
Str & \texttt{Value::Str(Rc<String>)} & \texttt{"hello"} & UTF-8 string, reference-counted \\
Array & \texttt{Value::Array(Rc<RefCell<Vec>>)} & \texttt{[1, 2, 3]} & Dynamic, mutable array \\
Hash & \texttt{Value::Hash(Rc<RefCell<HashMap>>)} & \texttt{\{name => "John"\}} & String-keyed dictionary \\
Object & \texttt{Value::Object} & \texttt{new Point(3, 4)} & OOP object with fields and methods \\
Class & \texttt{Value::Class} & \texttt{class Point \{ ... \}} & Class definition (stored value) \\
Ref & \texttt{Value::Ref(Rc<RefCell<Value>>)} & \texttt{\textbackslash expr} & Shared mutable reference \\
Function & \texttt{Value::NativeFn} or \texttt{Value::UserFn} & \texttt{fn(x) \{ ... \}} & Callable functions \\
\bottomrule
\end{tabular}
\end{table}

\subsection{Variable Declaration and Assignment}

Variables are declared with \texttt{let} and can be reassigned with \texttt{=}:

\begin{tcolorbox}[colback=codebg, colframe=darkpurple, colbacktitle=darkpurple, title=\textbf{\textcolor{codegold}{Variable Declaration and Assignment}}, fonttitle=\bfseries, sharp corners]
\begin{lstlisting}
# Declaration
let name = "CRISP";
let count = 0;
let @numbers = [1, 2, 3];

# Assignment (reassign existing variable)
count = count + 1;
name = "New Name";

# Assignment to undeclared variable is an error
# x = 5;  # Error: Undefined variable: x

# Constants
const PI = 3.14159;
# PI = 3.0;  # Would be an error (immutable)
\end{lstlisting}
\end{tcolorbox}

\subsection{Sigils}

CRISP supports Perl-inspired optional sigils for type annotation:

\begin{table}[H]
\centering
\rowcolors{2}{softviolet!30}{white}
\caption{\textcolor{darkpurple}{CRISP Sigils}}
\vspace{0.3em}
\begin{tabular}{@{}lllp{5cm}@{}}
\toprule
\rowcolor{softvioletdark}
\textbf{\textcolor{white}{Sigil}} & \textbf{\textcolor{white}{Name}} & \textbf{\textcolor{white}{Example}} & \textbf{\textcolor{white}{Description}} \\
\midrule
\texttt{\$} & Scalar & \texttt{\$name}, \texttt{\$count} & Single value (string, number, bool) \\
\texttt{@} & Array & \texttt{@items}, \texttt{@numbers} & Ordered list \\
\texttt{\%} & Hash & \texttt{\%config}, \texttt{\%person} & Key-value dictionary \\
\texttt{\&} & Reference & \texttt{\&callback}, \texttt{\&fn\_ref} & Reference to a value \\
\bottomrule
\end{tabular}
\end{table}

Sigils are optional — CRISP infers the type from the value:

\begin{tcolorbox}[colback=codebg, colframe=darkpurple, colbacktitle=darkpurple, title=\textbf{\textcolor{codegold}{Sigil Usage}}, fonttitle=\bfseries, sharp corners]
\begin{lstlisting}
# With sigils (Perl style)
let $name = "CRISP";
let @fruits = ["apple", "banana"];
let %person = { name => "John", age => 30 };
let &callback = |x| => x * 2;

# Without sigils (modern style)
let name = "CRISP";
let fruits = ["apple", "banana"];
let person = { name => "John", age => 30 };
let callback = |x| => x * 2;
\end{lstlisting}
\end{tcolorbox}

\subsection{Type Introspection and Conversion}

\begin{tcolorbox}[colback=codebg, colframe=darkpurple, colbacktitle=darkpurple, title=\textbf{\textcolor{codegold}{Type Functions}}, fonttitle=\bfseries, sharp corners]
\begin{lstlisting}
# Get the type of a value
say type(42);            # "int"
say type("hello");       # "string"
say type([1, 2, 3]);    # "array"
say type({ a => 1 });   # "hash"
say type(null);          # "null"
say type(true);          # "bool"

# Parse string to number
let n = int("42");       # Int(42)
let f = float("3.14");   # Float(3.14)

# Invalid parsing returns null
let bad = int("hello");  # null
if bad == null {
    say "Not a number!";
}

# Length of collections
let len = len("CRISP");         # 5
let len2 = len([1, 2, 3]);     # 3
let len3 = len({ a => 1, b => 2 });  # 2
\end{lstlisting}
\end{tcolorbox}
